\subsection{Learned Baselines Trained on the Bank's Own Simulations}
\label{sec:fair-baselines}
The main comparison uses two previously fitted ExtraTrees adapters at their
original operating points. Here we ask separately how learned baselines do
when trained on the bank's simulation cohort. Three learners were fitted
on its 24 training records (72 events). We fixed thresholds on 24
calibration records by requiring zero alarms in normal operation, then
scored each learner once on 24 held-out records with the same event
association rule. At each decision the inputs measure changes in
16 phasor channels over the previous 0.5 and 1.5\,s relative
to a reference 1.5--4\,s earlier, whitened by that record's startup noise.
One learner receives these changes directly. Two receive their projections
onto the 18 catalog signatures, giving them the bank's static physical
signatures. A multinomial logistic model on those projections is the linear
control. Physical methods had already opened the held-out records, but
those records were not used to fit or tune any learner.

Table~\ref{tab:fair-baselines} reports the outcomes. The best learner
locates \FBBestTop{} of \FBEvents{} events at the source, compared with
\FBBankTop{} for the dynamic bank. The lowest learner count is \FBWorstTop{}. The
signature projections help most for line outages (\FBFisLines{} of 36)
but less for loads (\FBFisLoads{} of 36), whose small electromechanical
responses are poorly represented by static signatures. The learners emit
\FBBestExtras{}--\FBWorstExtras{} alarms outside scheduled events, versus
\FBBankExtras{} for the bank. All had zero alarms on held-out normal
records. The largest learner deficit is load attribution.
\begin{table}[t]\centering\scriptsize\setlength{\tabcolsep}{2pt}
\caption{Learned baselines trained on the simulated cohort of the bank, scored once on its held-out 72 events (36 lines, 36 loads) with the same event-association rule. Thresholds were set on calibration records for zero alarms in normal records. Top-1 counts the events detected with the correct source. The 72 events come from 12 independent records of six events each, so the last column counts the records in which all six were located (family and source), with its exact 95\% interval; the events of a record share one noise realization and operating point. Extras are alarms matching no event.}
\label{tab:fair-baselines}
\begin{tabular}{lcccccc}\toprule
Method & Detected & Top-1 & Lines & Loads & Extras & Records (95\%) \\\midrule
Dynamic bank & 72/72 & 72 & 36/36 & 36/36 & 7 & 12/12 (74--100) \\
Static physical & 48/72 & 40 & 32/36 & 8/36 & 25 & 0/12 (0--26) \\
Trees, phasor changes & 61/72 & 41 & 25/36 & 16/36 & 82 & 1/12 (0--38) \\
Trees, signatures & 58/72 & 49 & 32/36 & 17/36 & 30 & 0/12 (0--26) \\
Logistic, signatures & 55/72 & 32 & 27/36 & 5/36 & 71 & 0/12 (0--26) \\
\bottomrule\end{tabular}\end{table}

